
\documentclass{article}
\usepackage{colortbl}
\usepackage{makecell}
\usepackage{multirow}
\usepackage{supertabular}

\begin{document}

\newcounter{utterance}

\centering \large Interaction Transcript for game `dond', experiment `coop\_en', episode 10 with reka{-}flash{-}3.1{-}t0.0.
\vspace{24pt}

{ \footnotesize  \setcounter{utterance}{1}
\setlength{\tabcolsep}{0pt}
\begin{supertabular}{c@{$\;$}|p{.15\linewidth}@{}p{.15\linewidth}p{.15\linewidth}p{.15\linewidth}p{.15\linewidth}p{.15\linewidth}}
    \# & $\;$A & \multicolumn{4}{c}{Game Master} & $\;\:$B\\
    \hline

    \theutterance \stepcounter{utterance}  
    & & \multicolumn{4}{p{0.6\linewidth}}{
        \cellcolor[rgb]{0.9,0.9,0.9}{
            \makecell[{{p{\linewidth}}}]{
                \texttt{\tiny{[P1$\langle$GM]}}
                \texttt{You are playing a cooperative negotiation game in which you have to agree on how to divide a set of items among you and another player.} \\
\\ 
\texttt{Rules:} \\
\texttt{(a) You and the other player are given a set of items. Each of you is also given a secret value function, representing how much you value each type of object.} \\
\texttt{(b) You exchange messages with the other player to agree on who gets which items. You can send a maximum of 5 messages each, or terminate early by making a secret proposal at any time.} \\
\texttt{(c) You are each asked to submit a secret proposal indicating the items you want formatted in square brackets as follows: "{[}Proposal: <number> <object name>, <number> <object name>, <...>{]}"} \\
\texttt{(d) If your proposals are complementary, i.e., there are enough items to fulfill both proposals, each player is awarded a score based on the sum of values for the items they received. Otherwise, both of you get zero points.} \\
\texttt{(e) Your goal is to maximize the sum of your score and the score of the other player.} \\
\\ 
\texttt{Let us start.} \\
\\ 
\texttt{The set of available items is:} \\
\\ 
\texttt{2 pliers, 1 chalk, 1 nail, 1 bucket} \\
\\ 
\texttt{Your secret value function is:} \\
\\ 
\texttt{plier: 2, chalk: 3, nail: 2, bucket: 1} \\
\\ 
\texttt{IMPORTANT: Your messages, unless it is the secret proposal, are directly transmitted to the other player, so do not include any response to the rules or text announcing your message. To make a secret proposal, use the indicated format. Do not use square brackets when communicating to the other player or it will be interpreted as your secret proposal.} \\
            }
        }
    }
    & \\ \\

    \theutterance \stepcounter{utterance}  
    & \multicolumn{4}{p{0.6\linewidth}}{
        \cellcolor[rgb]{0.9,0.9,0.9}{
            \makecell[{{p{\linewidth}}}]{
                \texttt{\tiny{[P1$\rangle$GM]}}
                \texttt{<reasoning>} \\
\texttt{Let's see. I need to divide these items between me and the other player to maximize the total score. The items are 2 pliers, 1 chalk, 1 nail, 1 bucket. My values are plier: 2, chalk: 3, nail: 2, bucket: 1. The other player has their own values, but I don't know what they are. So I have to make a proposal that's fair in some way, but since I don't know their preferences, maybe I should aim for a division that gives me the maximum possible total value, assuming the other player would also want to maximize the total.} \\
\\ 
\texttt{First, let's list out all the items and their quantities. There are two pliers, which are duplicates. The other items are unique. So total items: 2 pliers, 1 chalk, 1 nail, 1 bucket. That's 5 items in total. Each player needs to get some subset of these, and the proposals have to be complementary, meaning together they cover all items without overlap. So if I propose taking one plier, the other player must take the other plier, and the rest of the items. Wait, but the proposals are from each player. So each player submits a proposal of what they want, and if their proposals together include all items exactly once, then it works. Otherwise, both get zero.} \\
\\ 
\texttt{So the goal is to make a proposal such that if the other player's proposal is complementary, the total of our individual scores is maximized. But since I don't know the other player's proposal, maybe I need to propose a division that is Pareto efficient, meaning that neither player can be made better off without making the other worse off. That way, it's a good compromise.} \\
\\ 
\texttt{Alternatively, maybe I should calculate the total value of all items and see how to split them. Let's compute the total value from my perspective. My values: 2 pliers (each worth 2), so 2*2=4. Chalk is 3, nail 2, bucket 1. Total is 4+3+2+1=10. But the other player's total would be different based on their values. But since I don't know theirs, maybe I should aim for a split that gives me a certain percentage of the total, but how?} \\
\\ 
\texttt{Alternatively, think about the items. The highest value items} \\
            }
        }
    }
    & & \\ \\

    \theutterance \stepcounter{utterance}  
    & & & \multicolumn{2}{p{0.3\linewidth}}{
        \cellcolor[rgb]{0.9,0.9,0.9}{
            \makecell[{{p{\linewidth}}}]{
                \texttt{\tiny{[GM$|$GM]}}
                \texttt{continue} \\
            }
        }
    }
    & & \\ \\

    \theutterance \stepcounter{utterance}  
    & & & \multicolumn{2}{p{0.3\linewidth}}{
        \cellcolor[rgb]{0.9,0.9,0.9}{
            \makecell[{{p{\linewidth}}}]{
                \texttt{\tiny{[GM$|$GM]}}
                \texttt{<reasoning>} \\
\texttt{Let's see. I need to divide these items between me and the other player to maximize the total score. The items are 2 pliers, 1 chalk, 1 nail, 1 bucket. My values are plier: 2, chalk: 3, nail: 2, bucket: 1. The other player has their own values, but I don't know what they are. So I have to make a proposal that's fair in some way, but since I don't know their preferences, maybe I should aim for a division that gives me the maximum possible total value, assuming the other player would also want to maximize the total.} \\
\\ 
\texttt{First, let's list out all the items and their quantities. There are two pliers, which are duplicates. The other items are unique. So total items: 2 pliers, 1 chalk, 1 nail, 1 bucket. That's 5 items in total. Each player needs to get some subset of these, and the proposals have to be complementary, meaning together they cover all items without overlap. So if I propose taking one plier, the other player must take the other plier, and the rest of the items. Wait, but the proposals are from each player. So each player submits a proposal of what they want, and if their proposals together include all items exactly once, then it works. Otherwise, both get zero.} \\
\\ 
\texttt{So the goal is to make a proposal such that if the other player's proposal is complementary, the total of our individual scores is maximized. But since I don't know the other player's proposal, maybe I need to propose a division that is Pareto efficient, meaning that neither player can be made better off without making the other worse off. That way, it's a good compromise.} \\
\\ 
\texttt{Alternatively, maybe I should calculate the total value of all items and see how to split them. Let's compute the total value from my perspective. My values: 2 pliers (each worth 2), so 2*2=4. Chalk is 3, nail 2, bucket 1. Total is 4+3+2+1=10. But the other player's total would be different based on their values. But since I don't know theirs, maybe I should aim for a split that gives me a certain percentage of the total, but how?} \\
\\ 
\texttt{Alternatively, think about the items. The highest value items} \\
            }
        }
    }
    & & \\ \\

    \theutterance \stepcounter{utterance}  
    & & \multicolumn{4}{p{0.6\linewidth}}{
        \cellcolor[rgb]{0.9,0.9,0.9}{
            \makecell[{{p{\linewidth}}}]{
                \texttt{\tiny{[GM$\rangle$P2]}}
                \texttt{You are playing a cooperative negotiation game in which you have to agree on how to divide a set of items among you and another player.} \\
\\ 
\texttt{Rules:} \\
\texttt{(a) You and the other player are given a set of items. Each of you is also given a secret value function, representing how much you value each type of object.} \\
\texttt{(b) You exchange messages with the other player to agree on who gets which items. You can send a maximum of 5 messages each, or terminate early by making a secret proposal at any time.} \\
\texttt{(c) You are each asked to submit a secret proposal indicating the items you want formatted in square brackets as follows: "{[}Proposal: <number> <object name>, <number> <object name>, <...>{]}"} \\
\texttt{(d) If your proposals are complementary, i.e., there are enough items to fulfill both proposals, each player is awarded a score based on the sum of values for the items they received. Otherwise, both of you get zero points.} \\
\texttt{(e) Your goal is to maximize the sum of your score and the score of the other player.} \\
\\ 
\texttt{Let us start.} \\
\\ 
\texttt{The set of available items is:} \\
\\ 
\texttt{2 pliers, 1 chalk, 1 nail, 1 bucket} \\
\\ 
\texttt{Your secret value function is:} \\
\\ 
\texttt{plier: 0, chalk: 2, nail: 5, bucket: 3} \\
\\ 
\texttt{IMPORTANT: Your messages, unless it is the secret proposal, are directly transmitted to the other player, so do not include any response to the rules or text announcing your message. To make a secret proposal, use the indicated format. Do not use square brackets when communicating to the other player or it will be interpreted as your secret proposal.} \\
\\ 
\\ 
\texttt{<reasoning>} \\
\texttt{Let's see. I need to divide these items between me and the other player to maximize the total score. The items are 2 pliers, 1 chalk, 1 nail, 1 bucket. My values are plier: 2, chalk: 3, nail: 2, bucket: 1. The other player has their own values, but I don't know what they are. So I have to make a proposal that's fair in some way, but since I don't know their preferences, maybe I should aim for a division that gives me the maximum possible total value, assuming the other player would also want to maximize the total.} \\
\\ 
\texttt{First, let's list out all the items and their quantities. There are two pliers, which are duplicates. The other items are unique. So total items: 2 pliers, 1 chalk, 1 nail, 1 bucket. That's 5 items in total. Each player needs to get some subset of these, and the proposals have to be complementary, meaning together they cover all items without overlap. So if I propose taking one plier, the other player must take the other plier, and the rest of the items. Wait, but the proposals are from each player. So each player submits a proposal of what they want, and if their proposals together include all items exactly once, then it works. Otherwise, both get zero.} \\
\\ 
\texttt{So the goal is to make a proposal such that if the other player's proposal is complementary, the total of our individual scores is maximized. But since I don't know the other player's proposal, maybe I need to propose a division that is Pareto efficient, meaning that neither player can be made better off without making the other worse off. That way, it's a good compromise.} \\
\\ 
\texttt{Alternatively, maybe I should calculate the total value of all items and see how to split them. Let's compute the total value from my perspective. My values: 2 pliers (each worth 2), so 2*2=4. Chalk is 3, nail 2, bucket 1. Total is 4+3+2+1=10. But the other player's total would be different based on their values. But since I don't know theirs, maybe I should aim for a split that gives me a certain percentage of the total, but how?} \\
\\ 
\texttt{Alternatively, think about the items. The highest value items} \\
            }
        }
    }
    & \\ \\

    \theutterance \stepcounter{utterance}  
    & & & \multicolumn{4}{p{0.6\linewidth}}{
        \cellcolor[rgb]{0.9,0.9,0.9}{
            \makecell[{{p{\linewidth}}}]{
                \texttt{\tiny{[GM$\langle$P2]}}
                \texttt{<reasoning>} \\
\texttt{Let me try to figure out how to approach this. The goal is to maximize the sum of both players' scores. Since I don't know the other player's value function, I need to make a proposal that's fair in a way that balances both our potential gains. Maybe I should look at the items and see which ones are more valuable to me and try to split them in a way that's mutually beneficial.} \\
\\ 
\texttt{First, let's list out the items and their quantities again: 2 pliers, 1 chalk, 1 nail, 1 bucket. My values are plier: 2 each, chalk: 3, nail: 2, bucket: 1. The other player's values are unknown, but they must have some. The total value from my side is 2*2 (pliers) + 3 (chalk) + 2 (nail) + 1 (bucket) = 4 + 3 + 2 + 1 = 10. But the other player's total would vary.} \\
\\ 
\texttt{Since the other player's values are secret, I can't directly optimize for their preferences. However, to maximize the total score, I should aim for a division where the sum of our individual scores is as high as possible. This might involve giving them items that are more valuable to them and keeping items that are more valuable to me, but without knowing their values, it's tricky.} \\
\\ 
\texttt{One approach is to propose a division that gives the other player the items I value less, and keep the ones I value more. But since I don't know their values, maybe I should focus on splitting the items in a way that the total value for both is maximized. Alternatively, use the concept of a Nash bargaining solution, which maximizes the product of the players' utilities. But without knowing their utilities, that's hard.} \\
\\ 
\texttt{Alternatively, think about the items. The highest value items for me are chalk (3) and nail (2). The pliers are each 2, and bucket is 1. The other player might value some items higher. Maybe I should try to keep the chalk and nail for myself, and offer the pliers and bucket. But there are two pliers. If I take both, that's 4 value, but maybe the other player values them more. Wait, but I have two pliers. If I take both, the other player gets nothing else. But maybe the other player} \\
            }
        }
    }
     \\ \\

    \theutterance \stepcounter{utterance}  
    & & & \multicolumn{2}{p{0.3\linewidth}}{
        \cellcolor[rgb]{0.9,0.9,0.9}{
            \makecell[{{p{\linewidth}}}]{
                \texttt{\tiny{[GM$|$GM]}}
                \texttt{continue} \\
            }
        }
    }
    & & \\ \\

    \theutterance \stepcounter{utterance}  
    & & & \multicolumn{2}{p{0.3\linewidth}}{
        \cellcolor[rgb]{0.9,0.9,0.9}{
            \makecell[{{p{\linewidth}}}]{
                \texttt{\tiny{[GM$|$GM]}}
                \texttt{<reasoning>} \\
\texttt{Let me try to figure out how to approach this. The goal is to maximize the sum of both players' scores. Since I don't know the other player's value function, I need to make a proposal that's fair in a way that balances both our potential gains. Maybe I should look at the items and see which ones are more valuable to me and try to split them in a way that's mutually beneficial.} \\
\\ 
\texttt{First, let's list out the items and their quantities again: 2 pliers, 1 chalk, 1 nail, 1 bucket. My values are plier: 2 each, chalk: 3, nail: 2, bucket: 1. The other player's values are unknown, but they must have some. The total value from my side is 2*2 (pliers) + 3 (chalk) + 2 (nail) + 1 (bucket) = 4 + 3 + 2 + 1 = 10. But the other player's total would vary.} \\
\\ 
\texttt{Since the other player's values are secret, I can't directly optimize for their preferences. However, to maximize the total score, I should aim for a division where the sum of our individual scores is as high as possible. This might involve giving them items that are more valuable to them and keeping items that are more valuable to me, but without knowing their values, it's tricky.} \\
\\ 
\texttt{One approach is to propose a division that gives the other player the items I value less, and keep the ones I value more. But since I don't know their values, maybe I should focus on splitting the items in a way that the total value for both is maximized. Alternatively, use the concept of a Nash bargaining solution, which maximizes the product of the players' utilities. But without knowing their utilities, that's hard.} \\
\\ 
\texttt{Alternatively, think about the items. The highest value items for me are chalk (3) and nail (2). The pliers are each 2, and bucket is 1. The other player might value some items higher. Maybe I should try to keep the chalk and nail for myself, and offer the pliers and bucket. But there are two pliers. If I take both, that's 4 value, but maybe the other player values them more. Wait, but I have two pliers. If I take both, the other player gets nothing else. But maybe the other player} \\
            }
        }
    }
    & & \\ \\

    \theutterance \stepcounter{utterance}  
    & & \multicolumn{4}{p{0.6\linewidth}}{
        \cellcolor[rgb]{0.9,0.9,0.9}{
            \makecell[{{p{\linewidth}}}]{
                \texttt{\tiny{[P1$\langle$GM]}}
                \texttt{<reasoning>} \\
\texttt{Let me try to figure out how to approach this. The goal is to maximize the sum of both players' scores. Since I don't know the other player's value function, I need to make a proposal that's fair in a way that balances both our potential gains. Maybe I should look at the items and see which ones are more valuable to me and try to split them in a way that's mutually beneficial.} \\
\\ 
\texttt{First, let's list out the items and their quantities again: 2 pliers, 1 chalk, 1 nail, 1 bucket. My values are plier: 2 each, chalk: 3, nail: 2, bucket: 1. The other player's values are unknown, but they must have some. The total value from my side is 2*2 (pliers) + 3 (chalk) + 2 (nail) + 1 (bucket) = 4 + 3 + 2 + 1 = 10. But the other player's total would vary.} \\
\\ 
\texttt{Since the other player's values are secret, I can't directly optimize for their preferences. However, to maximize the total score, I should aim for a division where the sum of our individual scores is as high as possible. This might involve giving them items that are more valuable to them and keeping items that are more valuable to me, but without knowing their values, it's tricky.} \\
\\ 
\texttt{One approach is to propose a division that gives the other player the items I value less, and keep the ones I value more. But since I don't know their values, maybe I should focus on splitting the items in a way that the total value for both is maximized. Alternatively, use the concept of a Nash bargaining solution, which maximizes the product of the players' utilities. But without knowing their utilities, that's hard.} \\
\\ 
\texttt{Alternatively, think about the items. The highest value items for me are chalk (3) and nail (2). The pliers are each 2, and bucket is 1. The other player might value some items higher. Maybe I should try to keep the chalk and nail for myself, and offer the pliers and bucket. But there are two pliers. If I take both, that's 4 value, but maybe the other player values them more. Wait, but I have two pliers. If I take both, the other player gets nothing else. But maybe the other player} \\
            }
        }
    }
    & \\ \\

    \theutterance \stepcounter{utterance}  
    & \multicolumn{4}{p{0.6\linewidth}}{
        \cellcolor[rgb]{0.9,0.9,0.9}{
            \makecell[{{p{\linewidth}}}]{
                \texttt{\tiny{[P1$\rangle$GM]}}
                \texttt{<reasoning>} \\
\texttt{Let me try to work through this step by step. The goal is to divide the items so that the sum of both our scores is maximized. Since I don't know the other player's values, I need to make a proposal that's as fair as possible, assuming they want to maximize the total too.} \\
\\ 
\texttt{First, let's list out the items and their quantities again to make sure I have them right: 2 pliers, 1 chalk, 1 nail, 1 bucket. My values are plier: 2 each, chalk: 3, nail: 2, bucket: 1. The other player's values are unknown, but they must have some. The total value from my perspective is 2*2 (pliers) + 3 (chalk) + 2 (nail) + 1 (bucket) = 4 + 3 + 2 + 1 = 10. But the other player's total would depend on their values.} \\
\\ 
\texttt{Since the other player's values are secret, I can't directly optimize for their preferences. However, to maximize the total score, I should aim for a division where the sum of our individual scores is as high as possible. This might involve giving them items that are more valuable to them and keeping the ones I value more, but without knowing their values, it's tricky.} \\
\\ 
\texttt{One approach is to propose a division that gives the other player the items I value less, and keep the ones I value more. But since I don't know their values, maybe I should focus on splitting the items in a way that the total value for both is maximized. Alternatively, use the concept of a Nash bargaining solution, which maximizes the product of the players' utilities. But without knowing their utilities, that's hard.} \\
\\ 
\texttt{Alternatively, think about the items. The highest value items for me are chalk (3) and nail (2). The pliers are each 2, and bucket is 1. The other player might value some items higher. Maybe I should try to keep the chalk and nail for myself, and offer the pliers and bucket. But there are two pliers. If I take both, that's 4 value, but maybe the other player values them more. Wait, but I have two pliers. If I take both, the other player gets nothing else. But maybe the other player values the chalk or nail more. Hmm.} \\
\\ 
\texttt{Alternatively, maybe split} \\
            }
        }
    }
    & & \\ \\

    \theutterance \stepcounter{utterance}  
    & & & \multicolumn{2}{p{0.3\linewidth}}{
        \cellcolor[rgb]{0.9,0.9,0.9}{
            \makecell[{{p{\linewidth}}}]{
                \texttt{\tiny{[GM$|$GM]}}
                \texttt{continue} \\
            }
        }
    }
    & & \\ \\

    \theutterance \stepcounter{utterance}  
    & & & \multicolumn{2}{p{0.3\linewidth}}{
        \cellcolor[rgb]{0.9,0.9,0.9}{
            \makecell[{{p{\linewidth}}}]{
                \texttt{\tiny{[GM$|$GM]}}
                \texttt{<reasoning>} \\
\texttt{Let me try to work through this step by step. The goal is to divide the items so that the sum of both our scores is maximized. Since I don't know the other player's values, I need to make a proposal that's as fair as possible, assuming they want to maximize the total too.} \\
\\ 
\texttt{First, let's list out the items and their quantities again to make sure I have them right: 2 pliers, 1 chalk, 1 nail, 1 bucket. My values are plier: 2 each, chalk: 3, nail: 2, bucket: 1. The other player's values are unknown, but they must have some. The total value from my perspective is 2*2 (pliers) + 3 (chalk) + 2 (nail) + 1 (bucket) = 4 + 3 + 2 + 1 = 10. But the other player's total would depend on their values.} \\
\\ 
\texttt{Since the other player's values are secret, I can't directly optimize for their preferences. However, to maximize the total score, I should aim for a division where the sum of our individual scores is as high as possible. This might involve giving them items that are more valuable to them and keeping the ones I value more, but without knowing their values, it's tricky.} \\
\\ 
\texttt{One approach is to propose a division that gives the other player the items I value less, and keep the ones I value more. But since I don't know their values, maybe I should focus on splitting the items in a way that the total value for both is maximized. Alternatively, use the concept of a Nash bargaining solution, which maximizes the product of the players' utilities. But without knowing their utilities, that's hard.} \\
\\ 
\texttt{Alternatively, think about the items. The highest value items for me are chalk (3) and nail (2). The pliers are each 2, and bucket is 1. The other player might value some items higher. Maybe I should try to keep the chalk and nail for myself, and offer the pliers and bucket. But there are two pliers. If I take both, that's 4 value, but maybe the other player values them more. Wait, but I have two pliers. If I take both, the other player gets nothing else. But maybe the other player values the chalk or nail more. Hmm.} \\
\\ 
\texttt{Alternatively, maybe split} \\
            }
        }
    }
    & & \\ \\

    \theutterance \stepcounter{utterance}  
    & & \multicolumn{4}{p{0.6\linewidth}}{
        \cellcolor[rgb]{0.9,0.9,0.9}{
            \makecell[{{p{\linewidth}}}]{
                \texttt{\tiny{[GM$\rangle$P2]}}
                \texttt{<reasoning>} \\
\texttt{Let me try to work through this step by step. The goal is to divide the items so that the sum of both our scores is maximized. Since I don't know the other player's values, I need to make a proposal that's as fair as possible, assuming they want to maximize the total too.} \\
\\ 
\texttt{First, let's list out the items and their quantities again to make sure I have them right: 2 pliers, 1 chalk, 1 nail, 1 bucket. My values are plier: 2 each, chalk: 3, nail: 2, bucket: 1. The other player's values are unknown, but they must have some. The total value from my perspective is 2*2 (pliers) + 3 (chalk) + 2 (nail) + 1 (bucket) = 4 + 3 + 2 + 1 = 10. But the other player's total would depend on their values.} \\
\\ 
\texttt{Since the other player's values are secret, I can't directly optimize for their preferences. However, to maximize the total score, I should aim for a division where the sum of our individual scores is as high as possible. This might involve giving them items that are more valuable to them and keeping the ones I value more, but without knowing their values, it's tricky.} \\
\\ 
\texttt{One approach is to propose a division that gives the other player the items I value less, and keep the ones I value more. But since I don't know their values, maybe I should focus on splitting the items in a way that the total value for both is maximized. Alternatively, use the concept of a Nash bargaining solution, which maximizes the product of the players' utilities. But without knowing their utilities, that's hard.} \\
\\ 
\texttt{Alternatively, think about the items. The highest value items for me are chalk (3) and nail (2). The pliers are each 2, and bucket is 1. The other player might value some items higher. Maybe I should try to keep the chalk and nail for myself, and offer the pliers and bucket. But there are two pliers. If I take both, that's 4 value, but maybe the other player values them more. Wait, but I have two pliers. If I take both, the other player gets nothing else. But maybe the other player values the chalk or nail more. Hmm.} \\
\\ 
\texttt{Alternatively, maybe split} \\
            }
        }
    }
    & \\ \\

    \theutterance \stepcounter{utterance}  
    & & & \multicolumn{4}{p{0.6\linewidth}}{
        \cellcolor[rgb]{0.9,0.9,0.9}{
            \makecell[{{p{\linewidth}}}]{
                \texttt{\tiny{[GM$\langle$P2]}}
                \texttt{<reasoning>} \\
\texttt{Let me try to figure out the best way to divide these items. The goal is to maximize the total score for both players. Since I don't know the other player's values, I need to make a proposal that's fair and likely to be accepted. Let's start by evaluating the items based on my values.} \\
\\ 
\texttt{My values:} \\
\texttt{{-} Pliers: 2 each (total 4 for both)} \\
\texttt{{-} Chalk: 3} \\
\texttt{{-} Nail: 2} \\
\texttt{{-} Bucket: 1} \\
\\ 
\texttt{Total value I assign: 4 + 3 + 2 + 1 = 10.} \\
\\ 
\texttt{But the other player's values are unknown. However, to maximize the total, I should aim for a division where the sum of both our scores is as high as possible. This might mean balancing the items so that neither of us gets too much or too little.} \\
\\ 
\texttt{Let's list all possible items and think about how to split them. The items are 2 pliers, 1 chalk, 1 nail, 1 bucket. Each player needs to get some subset, and the proposals must be complementary.} \\
\\ 
\texttt{If I take both pliers, that's 4 value for me. Then the other player gets chalk, nail, bucket. Their value for those would depend on their own values. But maybe that's too much for me. Alternatively, split the pliers: give one to the other player and keep one. Then I have one plier (2), chalk (3), nail (2), bucket (1) for a total of 8. The other player gets one plier (2) and whatever's left. Wait, no—if I take one plier, chalk, nail, bucket, that's 2+3+2+1=8. The other player gets the other plier (2). So their total would be 2. That's a big difference. Not good for total.} \\
\\ 
\texttt{Alternatively, maybe split the high{-}value items. Keep chalk (3) and nail (2) for myself, and split the pliers and bucket. So I take chalk, nail, and one plier (total 3+2+2=7). The other player gets the other plier (2) and bucket (1). Their total would be 3. Total combined is 10. That's good. But maybe the other player values the bucket more? Not sure.} \\
\\ 
\texttt{Wait, but the bucket is only worth 1 to} \\
            }
        }
    }
     \\ \\

    \theutterance \stepcounter{utterance}  
    & & & \multicolumn{2}{p{0.3\linewidth}}{
        \cellcolor[rgb]{0.9,0.9,0.9}{
            \makecell[{{p{\linewidth}}}]{
                \texttt{\tiny{[GM$|$GM]}}
                \texttt{continue} \\
            }
        }
    }
    & & \\ \\

    \theutterance \stepcounter{utterance}  
    & & & \multicolumn{2}{p{0.3\linewidth}}{
        \cellcolor[rgb]{0.9,0.9,0.9}{
            \makecell[{{p{\linewidth}}}]{
                \texttt{\tiny{[GM$|$GM]}}
                \texttt{<reasoning>} \\
\texttt{Let me try to figure out the best way to divide these items. The goal is to maximize the total score for both players. Since I don't know the other player's values, I need to make a proposal that's fair and likely to be accepted. Let's start by evaluating the items based on my values.} \\
\\ 
\texttt{My values:} \\
\texttt{{-} Pliers: 2 each (total 4 for both)} \\
\texttt{{-} Chalk: 3} \\
\texttt{{-} Nail: 2} \\
\texttt{{-} Bucket: 1} \\
\\ 
\texttt{Total value I assign: 4 + 3 + 2 + 1 = 10.} \\
\\ 
\texttt{But the other player's values are unknown. However, to maximize the total, I should aim for a division where the sum of both our scores is as high as possible. This might mean balancing the items so that neither of us gets too much or too little.} \\
\\ 
\texttt{Let's list all possible items and think about how to split them. The items are 2 pliers, 1 chalk, 1 nail, 1 bucket. Each player needs to get some subset, and the proposals must be complementary.} \\
\\ 
\texttt{If I take both pliers, that's 4 value for me. Then the other player gets chalk, nail, bucket. Their value for those would depend on their own values. But maybe that's too much for me. Alternatively, split the pliers: give one to the other player and keep one. Then I have one plier (2), chalk (3), nail (2), bucket (1) for a total of 8. The other player gets one plier (2) and whatever's left. Wait, no—if I take one plier, chalk, nail, bucket, that's 2+3+2+1=8. The other player gets the other plier (2). So their total would be 2. That's a big difference. Not good for total.} \\
\\ 
\texttt{Alternatively, maybe split the high{-}value items. Keep chalk (3) and nail (2) for myself, and split the pliers and bucket. So I take chalk, nail, and one plier (total 3+2+2=7). The other player gets the other plier (2) and bucket (1). Their total would be 3. Total combined is 10. That's good. But maybe the other player values the bucket more? Not sure.} \\
\\ 
\texttt{Wait, but the bucket is only worth 1 to} \\
            }
        }
    }
    & & \\ \\

    \theutterance \stepcounter{utterance}  
    & & \multicolumn{4}{p{0.6\linewidth}}{
        \cellcolor[rgb]{0.9,0.9,0.9}{
            \makecell[{{p{\linewidth}}}]{
                \texttt{\tiny{[P1$\langle$GM]}}
                \texttt{<reasoning>} \\
\texttt{Let me try to figure out the best way to divide these items. The goal is to maximize the total score for both players. Since I don't know the other player's values, I need to make a proposal that's fair and likely to be accepted. Let's start by evaluating the items based on my values.} \\
\\ 
\texttt{My values:} \\
\texttt{{-} Pliers: 2 each (total 4 for both)} \\
\texttt{{-} Chalk: 3} \\
\texttt{{-} Nail: 2} \\
\texttt{{-} Bucket: 1} \\
\\ 
\texttt{Total value I assign: 4 + 3 + 2 + 1 = 10.} \\
\\ 
\texttt{But the other player's values are unknown. However, to maximize the total, I should aim for a division where the sum of both our scores is as high as possible. This might mean balancing the items so that neither of us gets too much or too little.} \\
\\ 
\texttt{Let's list all possible items and think about how to split them. The items are 2 pliers, 1 chalk, 1 nail, 1 bucket. Each player needs to get some subset, and the proposals must be complementary.} \\
\\ 
\texttt{If I take both pliers, that's 4 value for me. Then the other player gets chalk, nail, bucket. Their value for those would depend on their own values. But maybe that's too much for me. Alternatively, split the pliers: give one to the other player and keep one. Then I have one plier (2), chalk (3), nail (2), bucket (1) for a total of 8. The other player gets one plier (2) and whatever's left. Wait, no—if I take one plier, chalk, nail, bucket, that's 2+3+2+1=8. The other player gets the other plier (2). So their total would be 2. That's a big difference. Not good for total.} \\
\\ 
\texttt{Alternatively, maybe split the high{-}value items. Keep chalk (3) and nail (2) for myself, and split the pliers and bucket. So I take chalk, nail, and one plier (total 3+2+2=7). The other player gets the other plier (2) and bucket (1). Their total would be 3. Total combined is 10. That's good. But maybe the other player values the bucket more? Not sure.} \\
\\ 
\texttt{Wait, but the bucket is only worth 1 to} \\
            }
        }
    }
    & \\ \\

    \theutterance \stepcounter{utterance}  
    & \multicolumn{4}{p{0.6\linewidth}}{
        \cellcolor[rgb]{0.9,0.9,0.9}{
            \makecell[{{p{\linewidth}}}]{
                \texttt{\tiny{[P1$\rangle$GM]}}
                \texttt{<reasoning>} \\
\texttt{Let's try to figure out the best way to divide these items. The goal is to maximize the total score for both players. Since I don't know the other player's values, I need to make a proposal that's fair and likely to be accepted. Let's start by evaluating the items based on my values.} \\
\\ 
\texttt{My values:} \\
\texttt{{-} Pliers: 2 each (total 4 for both)} \\
\texttt{{-} Chalk: 3} \\
\texttt{{-} Nail: 2} \\
\texttt{{-} Bucket: 1} \\
\\ 
\texttt{Total value I assign: 4 + 3 + 2 + 1 = 10.} \\
\\ 
\texttt{But the other player's values are unknown. However, to maximize the total, I should aim for a division where the sum of both our scores is as high as possible. This might mean balancing the items so that neither of us gets too much or too little.} \\
\\ 
\texttt{Let's list all possible items and think about how to split them. The items are 2 pliers, 1 chalk, 1 nail, 1 bucket. Each player needs to get some subset, and the proposals must be complementary.} \\
\\ 
\texttt{If I take both pliers, that's 4 value for me. Then the other player gets chalk, nail, bucket. Their value for those would depend on their own values. But maybe that's too much for me. Alternatively, split the pliers: give one to the other player and keep one. Then I have one plier (2), chalk (3), nail (2), bucket (1) for a total of 8. The other player gets one plier (2) and whatever's left. Wait, no—if I take one plier, chalk, nail, bucket, that's 2+3+2+1=8. The other player gets the other plier (2). So their total would be 2. That's a big difference. Not good for total.} \\
\\ 
\texttt{Alternatively, maybe split the high{-}value items. Keep chalk (3) and nail (2) for myself, and split the pliers and bucket. So I take chalk, nail, and one plier (total 3+2+2=7). The other player gets the other plier (2) and bucket (1). Their total would be 3. Total combined is 10. That's good. But maybe the other player values the bucket more? Not sure.} \\
\\ 
\texttt{Wait, but the bucket is only worth 1 to} \\
            }
        }
    }
    & & \\ \\

    \theutterance \stepcounter{utterance}  
    & & & \multicolumn{2}{p{0.3\linewidth}}{
        \cellcolor[rgb]{0.9,0.9,0.9}{
            \makecell[{{p{\linewidth}}}]{
                \texttt{\tiny{[GM$|$GM]}}
                \texttt{continue} \\
            }
        }
    }
    & & \\ \\

    \theutterance \stepcounter{utterance}  
    & & & \multicolumn{2}{p{0.3\linewidth}}{
        \cellcolor[rgb]{0.9,0.9,0.9}{
            \makecell[{{p{\linewidth}}}]{
                \texttt{\tiny{[GM$|$GM]}}
                \texttt{<reasoning>} \\
\texttt{Let's try to figure out the best way to divide these items. The goal is to maximize the total score for both players. Since I don't know the other player's values, I need to make a proposal that's fair and likely to be accepted. Let's start by evaluating the items based on my values.} \\
\\ 
\texttt{My values:} \\
\texttt{{-} Pliers: 2 each (total 4 for both)} \\
\texttt{{-} Chalk: 3} \\
\texttt{{-} Nail: 2} \\
\texttt{{-} Bucket: 1} \\
\\ 
\texttt{Total value I assign: 4 + 3 + 2 + 1 = 10.} \\
\\ 
\texttt{But the other player's values are unknown. However, to maximize the total, I should aim for a division where the sum of both our scores is as high as possible. This might mean balancing the items so that neither of us gets too much or too little.} \\
\\ 
\texttt{Let's list all possible items and think about how to split them. The items are 2 pliers, 1 chalk, 1 nail, 1 bucket. Each player needs to get some subset, and the proposals must be complementary.} \\
\\ 
\texttt{If I take both pliers, that's 4 value for me. Then the other player gets chalk, nail, bucket. Their value for those would depend on their own values. But maybe that's too much for me. Alternatively, split the pliers: give one to the other player and keep one. Then I have one plier (2), chalk (3), nail (2), bucket (1) for a total of 8. The other player gets one plier (2) and whatever's left. Wait, no—if I take one plier, chalk, nail, bucket, that's 2+3+2+1=8. The other player gets the other plier (2). So their total would be 2. That's a big difference. Not good for total.} \\
\\ 
\texttt{Alternatively, maybe split the high{-}value items. Keep chalk (3) and nail (2) for myself, and split the pliers and bucket. So I take chalk, nail, and one plier (total 3+2+2=7). The other player gets the other plier (2) and bucket (1). Their total would be 3. Total combined is 10. That's good. But maybe the other player values the bucket more? Not sure.} \\
\\ 
\texttt{Wait, but the bucket is only worth 1 to} \\
            }
        }
    }
    & & \\ \\

    \theutterance \stepcounter{utterance}  
    & & \multicolumn{4}{p{0.6\linewidth}}{
        \cellcolor[rgb]{0.9,0.9,0.9}{
            \makecell[{{p{\linewidth}}}]{
                \texttt{\tiny{[GM$\rangle$P2]}}
                \texttt{<reasoning>} \\
\texttt{Let's try to figure out the best way to divide these items. The goal is to maximize the total score for both players. Since I don't know the other player's values, I need to make a proposal that's fair and likely to be accepted. Let's start by evaluating the items based on my values.} \\
\\ 
\texttt{My values:} \\
\texttt{{-} Pliers: 2 each (total 4 for both)} \\
\texttt{{-} Chalk: 3} \\
\texttt{{-} Nail: 2} \\
\texttt{{-} Bucket: 1} \\
\\ 
\texttt{Total value I assign: 4 + 3 + 2 + 1 = 10.} \\
\\ 
\texttt{But the other player's values are unknown. However, to maximize the total, I should aim for a division where the sum of both our scores is as high as possible. This might mean balancing the items so that neither of us gets too much or too little.} \\
\\ 
\texttt{Let's list all possible items and think about how to split them. The items are 2 pliers, 1 chalk, 1 nail, 1 bucket. Each player needs to get some subset, and the proposals must be complementary.} \\
\\ 
\texttt{If I take both pliers, that's 4 value for me. Then the other player gets chalk, nail, bucket. Their value for those would depend on their own values. But maybe that's too much for me. Alternatively, split the pliers: give one to the other player and keep one. Then I have one plier (2), chalk (3), nail (2), bucket (1) for a total of 8. The other player gets one plier (2) and whatever's left. Wait, no—if I take one plier, chalk, nail, bucket, that's 2+3+2+1=8. The other player gets the other plier (2). So their total would be 2. That's a big difference. Not good for total.} \\
\\ 
\texttt{Alternatively, maybe split the high{-}value items. Keep chalk (3) and nail (2) for myself, and split the pliers and bucket. So I take chalk, nail, and one plier (total 3+2+2=7). The other player gets the other plier (2) and bucket (1). Their total would be 3. Total combined is 10. That's good. But maybe the other player values the bucket more? Not sure.} \\
\\ 
\texttt{Wait, but the bucket is only worth 1 to} \\
            }
        }
    }
    & \\ \\

    \theutterance \stepcounter{utterance}  
    & & & \multicolumn{4}{p{0.6\linewidth}}{
        \cellcolor[rgb]{0.9,0.9,0.9}{
            \makecell[{{p{\linewidth}}}]{
                \texttt{\tiny{[GM$\langle$P2]}}
                \texttt{<reasoning>} \\
\texttt{Let me try to work through this step by step. The goal is to divide the items so that the sum of both our scores is maximized. Since I don't know the other player's values, I need to make a proposal that's fair and likely to be accepted. Let's start by evaluating the items based on my values.} \\
\\ 
\texttt{My values:} \\
\texttt{{-} Pliers: 2 each (total 4 for both)} \\
\texttt{{-} Chalk: 3} \\
\texttt{{-} Nail: 2} \\
\texttt{{-} Bucket: 1} \\
\\ 
\texttt{Total value I assign: 4 + 3 + 2 + 1 = 10.} \\
\\ 
\texttt{But the other player's values are unknown. However, to maximize the total, I should aim for a division where the sum of both our scores is as high as possible. This might mean balancing the items so that neither of us gets too much or too little.} \\
\\ 
\texttt{Let's list all possible items and think about how to split them. The items are 2 pliers, 1 chalk, 1 nail, 1 bucket. Each player needs to get some subset, and the proposals must be complementary.} \\
\\ 
\texttt{If I take both pliers, that's 4 value for me. Then the other player gets chalk, nail, bucket. Their value for those would depend on their own values. But maybe that's too much for me. Alternatively, split the pliers: give one to the other player and keep one. Then I have one plier (2), chalk (3), nail (2), bucket (1) for a total of 8. The other player gets one plier (2) and whatever's left. Wait, no—if I take one plier, chalk, nail, bucket, that's 2+3+2+1=8. The other player gets the other plier (2). So their total would be 2. That's a big difference. Not good for total.} \\
\\ 
\texttt{Alternatively, maybe split the high{-}value items. Keep chalk (3) and nail (2) for myself, and split the pliers and bucket. So I take chalk, nail, and one plier (total 3+2+2=7). The other player gets the other plier (2) and bucket (1). Their total would be 3. Total combined is 10. That's good. But maybe the other player values the bucket more? Not sure.} \\
\\ 
\texttt{Wait, but the bucket is only} \\
            }
        }
    }
     \\ \\

    \theutterance \stepcounter{utterance}  
    & & & \multicolumn{2}{p{0.3\linewidth}}{
        \cellcolor[rgb]{0.9,0.9,0.9}{
            \makecell[{{p{\linewidth}}}]{
                \texttt{\tiny{[GM$|$GM]}}
                \texttt{continue} \\
            }
        }
    }
    & & \\ \\

    \theutterance \stepcounter{utterance}  
    & & & \multicolumn{2}{p{0.3\linewidth}}{
        \cellcolor[rgb]{0.9,0.9,0.9}{
            \makecell[{{p{\linewidth}}}]{
                \texttt{\tiny{[GM$|$GM]}}
                \texttt{<reasoning>} \\
\texttt{Let me try to work through this step by step. The goal is to divide the items so that the sum of both our scores is maximized. Since I don't know the other player's values, I need to make a proposal that's fair and likely to be accepted. Let's start by evaluating the items based on my values.} \\
\\ 
\texttt{My values:} \\
\texttt{{-} Pliers: 2 each (total 4 for both)} \\
\texttt{{-} Chalk: 3} \\
\texttt{{-} Nail: 2} \\
\texttt{{-} Bucket: 1} \\
\\ 
\texttt{Total value I assign: 4 + 3 + 2 + 1 = 10.} \\
\\ 
\texttt{But the other player's values are unknown. However, to maximize the total, I should aim for a division where the sum of both our scores is as high as possible. This might mean balancing the items so that neither of us gets too much or too little.} \\
\\ 
\texttt{Let's list all possible items and think about how to split them. The items are 2 pliers, 1 chalk, 1 nail, 1 bucket. Each player needs to get some subset, and the proposals must be complementary.} \\
\\ 
\texttt{If I take both pliers, that's 4 value for me. Then the other player gets chalk, nail, bucket. Their value for those would depend on their own values. But maybe that's too much for me. Alternatively, split the pliers: give one to the other player and keep one. Then I have one plier (2), chalk (3), nail (2), bucket (1) for a total of 8. The other player gets one plier (2) and whatever's left. Wait, no—if I take one plier, chalk, nail, bucket, that's 2+3+2+1=8. The other player gets the other plier (2). So their total would be 2. That's a big difference. Not good for total.} \\
\\ 
\texttt{Alternatively, maybe split the high{-}value items. Keep chalk (3) and nail (2) for myself, and split the pliers and bucket. So I take chalk, nail, and one plier (total 3+2+2=7). The other player gets the other plier (2) and bucket (1). Their total would be 3. Total combined is 10. That's good. But maybe the other player values the bucket more? Not sure.} \\
\\ 
\texttt{Wait, but the bucket is only} \\
            }
        }
    }
    & & \\ \\

    \theutterance \stepcounter{utterance}  
    & & \multicolumn{4}{p{0.6\linewidth}}{
        \cellcolor[rgb]{0.9,0.9,0.9}{
            \makecell[{{p{\linewidth}}}]{
                \texttt{\tiny{[P1$\langle$GM]}}
                \texttt{<reasoning>} \\
\texttt{Let me try to work through this step by step. The goal is to divide the items so that the sum of both our scores is maximized. Since I don't know the other player's values, I need to make a proposal that's fair and likely to be accepted. Let's start by evaluating the items based on my values.} \\
\\ 
\texttt{My values:} \\
\texttt{{-} Pliers: 2 each (total 4 for both)} \\
\texttt{{-} Chalk: 3} \\
\texttt{{-} Nail: 2} \\
\texttt{{-} Bucket: 1} \\
\\ 
\texttt{Total value I assign: 4 + 3 + 2 + 1 = 10.} \\
\\ 
\texttt{But the other player's values are unknown. However, to maximize the total, I should aim for a division where the sum of both our scores is as high as possible. This might mean balancing the items so that neither of us gets too much or too little.} \\
\\ 
\texttt{Let's list all possible items and think about how to split them. The items are 2 pliers, 1 chalk, 1 nail, 1 bucket. Each player needs to get some subset, and the proposals must be complementary.} \\
\\ 
\texttt{If I take both pliers, that's 4 value for me. Then the other player gets chalk, nail, bucket. Their value for those would depend on their own values. But maybe that's too much for me. Alternatively, split the pliers: give one to the other player and keep one. Then I have one plier (2), chalk (3), nail (2), bucket (1) for a total of 8. The other player gets one plier (2) and whatever's left. Wait, no—if I take one plier, chalk, nail, bucket, that's 2+3+2+1=8. The other player gets the other plier (2). So their total would be 2. That's a big difference. Not good for total.} \\
\\ 
\texttt{Alternatively, maybe split the high{-}value items. Keep chalk (3) and nail (2) for myself, and split the pliers and bucket. So I take chalk, nail, and one plier (total 3+2+2=7). The other player gets the other plier (2) and bucket (1). Their total would be 3. Total combined is 10. That's good. But maybe the other player values the bucket more? Not sure.} \\
\\ 
\texttt{Wait, but the bucket is only} \\
            }
        }
    }
    & \\ \\

    \theutterance \stepcounter{utterance}  
    & \multicolumn{4}{p{0.6\linewidth}}{
        \cellcolor[rgb]{0.9,0.9,0.9}{
            \makecell[{{p{\linewidth}}}]{
                \texttt{\tiny{[P1$\rangle$GM]}}
                \texttt{<reasoning>} \\
\texttt{Let me try to figure out the best way to divide these items. The goal is to maximize the total score for both players. Since I don't know the other player's values, I need to make a proposal that's fair and likely to be accepted. Let's start by evaluating the items based on my values.} \\
\\ 
\texttt{My values:} \\
\texttt{{-} Pliers: 2 each (total 4 for both)} \\
\texttt{{-} Chalk: 3} \\
\texttt{{-} Nail: 2} \\
\texttt{{-} Bucket: 1} \\
\\ 
\texttt{Total value I assign: 4 + 3 + 2 + 1 = 10.} \\
\\ 
\texttt{But the other player's values are unknown. However, to maximize the total, I should aim for a division where the sum of both our scores is as high as possible. This might mean balancing the items so that neither of us gets too much or too little.} \\
\\ 
\texttt{Let's list all possible items and think about how to split them. The items are 2 pliers, 1 chalk, 1 nail, 1 bucket. Each player needs to get some subset, and the proposals must be complementary.} \\
\\ 
\texttt{If I take both pliers, that's 4 value for me. Then the other player gets chalk, nail, bucket. Their value for those would depend on their own values. But maybe that's too much for me. Alternatively, split the pliers: give one to the other player and keep one. Then I have one plier (2), chalk (3), nail (2), bucket (1) for a total of 8. The other player gets one plier (2) and whatever's left. Wait, no—if I take one plier, chalk, nail, bucket, that's 2+3+2+1=8. The other player gets the other plier (2). So their total would be 2. That's a big difference. Not good for total.} \\
\\ 
\texttt{Alternatively, maybe split the high{-}value items. Keep chalk (3) and nail (2) for myself, and split the pliers and bucket. So I take chalk, nail, and one plier (total 3+2+2=7). The other player gets the other plier (2) and bucket (1). Their total would be 3. Total combined is 10. That's good. But maybe the other player values the bucket more? Not sure.} \\
\\ 
\texttt{Wait, but the bucket is only worth 1 to} \\
            }
        }
    }
    & & \\ \\

    \theutterance \stepcounter{utterance}  
    & & & \multicolumn{2}{p{0.3\linewidth}}{
        \cellcolor[rgb]{0.9,0.9,0.9}{
            \makecell[{{p{\linewidth}}}]{
                \texttt{\tiny{[GM$|$GM]}}
                \texttt{continue} \\
            }
        }
    }
    & & \\ \\

    \theutterance \stepcounter{utterance}  
    & & & \multicolumn{2}{p{0.3\linewidth}}{
        \cellcolor[rgb]{0.9,0.9,0.9}{
            \makecell[{{p{\linewidth}}}]{
                \texttt{\tiny{[GM$|$GM]}}
                \texttt{<reasoning>} \\
\texttt{Let me try to figure out the best way to divide these items. The goal is to maximize the total score for both players. Since I don't know the other player's values, I need to make a proposal that's fair and likely to be accepted. Let's start by evaluating the items based on my values.} \\
\\ 
\texttt{My values:} \\
\texttt{{-} Pliers: 2 each (total 4 for both)} \\
\texttt{{-} Chalk: 3} \\
\texttt{{-} Nail: 2} \\
\texttt{{-} Bucket: 1} \\
\\ 
\texttt{Total value I assign: 4 + 3 + 2 + 1 = 10.} \\
\\ 
\texttt{But the other player's values are unknown. However, to maximize the total, I should aim for a division where the sum of both our scores is as high as possible. This might mean balancing the items so that neither of us gets too much or too little.} \\
\\ 
\texttt{Let's list all possible items and think about how to split them. The items are 2 pliers, 1 chalk, 1 nail, 1 bucket. Each player needs to get some subset, and the proposals must be complementary.} \\
\\ 
\texttt{If I take both pliers, that's 4 value for me. Then the other player gets chalk, nail, bucket. Their value for those would depend on their own values. But maybe that's too much for me. Alternatively, split the pliers: give one to the other player and keep one. Then I have one plier (2), chalk (3), nail (2), bucket (1) for a total of 8. The other player gets one plier (2) and whatever's left. Wait, no—if I take one plier, chalk, nail, bucket, that's 2+3+2+1=8. The other player gets the other plier (2). So their total would be 2. That's a big difference. Not good for total.} \\
\\ 
\texttt{Alternatively, maybe split the high{-}value items. Keep chalk (3) and nail (2) for myself, and split the pliers and bucket. So I take chalk, nail, and one plier (total 3+2+2=7). The other player gets the other plier (2) and bucket (1). Their total would be 3. Total combined is 10. That's good. But maybe the other player values the bucket more? Not sure.} \\
\\ 
\texttt{Wait, but the bucket is only worth 1 to} \\
            }
        }
    }
    & & \\ \\

    \theutterance \stepcounter{utterance}  
    & & \multicolumn{4}{p{0.6\linewidth}}{
        \cellcolor[rgb]{0.9,0.9,0.9}{
            \makecell[{{p{\linewidth}}}]{
                \texttt{\tiny{[GM$\rangle$P2]}}
                \texttt{<reasoning>} \\
\texttt{Let me try to figure out the best way to divide these items. The goal is to maximize the total score for both players. Since I don't know the other player's values, I need to make a proposal that's fair and likely to be accepted. Let's start by evaluating the items based on my values.} \\
\\ 
\texttt{My values:} \\
\texttt{{-} Pliers: 2 each (total 4 for both)} \\
\texttt{{-} Chalk: 3} \\
\texttt{{-} Nail: 2} \\
\texttt{{-} Bucket: 1} \\
\\ 
\texttt{Total value I assign: 4 + 3 + 2 + 1 = 10.} \\
\\ 
\texttt{But the other player's values are unknown. However, to maximize the total, I should aim for a division where the sum of both our scores is as high as possible. This might mean balancing the items so that neither of us gets too much or too little.} \\
\\ 
\texttt{Let's list all possible items and think about how to split them. The items are 2 pliers, 1 chalk, 1 nail, 1 bucket. Each player needs to get some subset, and the proposals must be complementary.} \\
\\ 
\texttt{If I take both pliers, that's 4 value for me. Then the other player gets chalk, nail, bucket. Their value for those would depend on their own values. But maybe that's too much for me. Alternatively, split the pliers: give one to the other player and keep one. Then I have one plier (2), chalk (3), nail (2), bucket (1) for a total of 8. The other player gets one plier (2) and whatever's left. Wait, no—if I take one plier, chalk, nail, bucket, that's 2+3+2+1=8. The other player gets the other plier (2). So their total would be 2. That's a big difference. Not good for total.} \\
\\ 
\texttt{Alternatively, maybe split the high{-}value items. Keep chalk (3) and nail (2) for myself, and split the pliers and bucket. So I take chalk, nail, and one plier (total 3+2+2=7). The other player gets the other plier (2) and bucket (1). Their total would be 3. Total combined is 10. That's good. But maybe the other player values the bucket more? Not sure.} \\
\\ 
\texttt{Wait, but the bucket is only worth 1 to} \\
            }
        }
    }
    & \\ \\

    \theutterance \stepcounter{utterance}  
    & & & \multicolumn{4}{p{0.6\linewidth}}{
        \cellcolor[rgb]{0.9,0.9,0.9}{
            \makecell[{{p{\linewidth}}}]{
                \texttt{\tiny{[GM$\langle$P2]}}
                \texttt{<reasoning>} \\
\texttt{Let me try to work through this step by step. The goal is to divide the items so that the sum of both our scores is maximized. Since I don't know the other player's values, I need to make a proposal that's fair and likely to be accepted. Let's start by evaluating the items based on my values.} \\
\\ 
\texttt{My values:} \\
\texttt{{-} Pliers: 2 each (total 4 for both)} \\
\texttt{{-} Chalk: 3} \\
\texttt{{-} Nail: 2} \\
\texttt{{-} Bucket: 1} \\
\\ 
\texttt{Total value I assign: 4 + 3 + 2 + 1 = 10.} \\
\\ 
\texttt{But the other player's values are unknown. However, to maximize the total, I should aim for a division where the sum of both our scores is as high as possible. This might mean balancing the items so that neither of us gets too much or too little.} \\
\\ 
\texttt{Let's list all possible items and think about how to split them. The items are 2 pliers, 1 chalk, 1 nail, 1 bucket. Each player needs to get some subset, and the proposals must be complementary.} \\
\\ 
\texttt{If I take both pliers, that's 4 value for me. Then the other player gets chalk, nail, bucket. Their value for those would depend on their own values. But maybe that's too much for me. Alternatively, split the pliers: give one to the other player and keep one. Then I have one plier (2), chalk (3), nail (2), bucket (1) for a total of 8. The other player gets one plier (2) and whatever's left. Wait, no—if I take one plier, chalk, nail, bucket, that's 2+3+2+1=8. The other player gets the other plier (2). So their total would be 2. That's a big difference. Not good for total.} \\
\\ 
\texttt{Alternatively, maybe split the high{-}value items. Keep chalk (3) and nail (2) for myself, and split the pliers and bucket. So I take chalk, nail, and one plier (total 3+2+2=7). The other player gets the other plier (2) and bucket (1). Their total would be 3. Total combined is 10. That's good. But maybe the other player values the bucket more? Not sure.} \\
\\ 
\texttt{Wait, but the bucket is only} \\
            }
        }
    }
     \\ \\

    \theutterance \stepcounter{utterance}  
    & & & \multicolumn{2}{p{0.3\linewidth}}{
        \cellcolor[rgb]{0.9,0.9,0.9}{
            \makecell[{{p{\linewidth}}}]{
                \texttt{\tiny{[GM$|$GM]}}
                \texttt{continue} \\
            }
        }
    }
    & & \\ \\

    \theutterance \stepcounter{utterance}  
    & & & \multicolumn{2}{p{0.3\linewidth}}{
        \cellcolor[rgb]{0.9,0.9,0.9}{
            \makecell[{{p{\linewidth}}}]{
                \texttt{\tiny{[GM$|$GM]}}
                \texttt{<reasoning>} \\
\texttt{Let me try to work through this step by step. The goal is to divide the items so that the sum of both our scores is maximized. Since I don't know the other player's values, I need to make a proposal that's fair and likely to be accepted. Let's start by evaluating the items based on my values.} \\
\\ 
\texttt{My values:} \\
\texttt{{-} Pliers: 2 each (total 4 for both)} \\
\texttt{{-} Chalk: 3} \\
\texttt{{-} Nail: 2} \\
\texttt{{-} Bucket: 1} \\
\\ 
\texttt{Total value I assign: 4 + 3 + 2 + 1 = 10.} \\
\\ 
\texttt{But the other player's values are unknown. However, to maximize the total, I should aim for a division where the sum of both our scores is as high as possible. This might mean balancing the items so that neither of us gets too much or too little.} \\
\\ 
\texttt{Let's list all possible items and think about how to split them. The items are 2 pliers, 1 chalk, 1 nail, 1 bucket. Each player needs to get some subset, and the proposals must be complementary.} \\
\\ 
\texttt{If I take both pliers, that's 4 value for me. Then the other player gets chalk, nail, bucket. Their value for those would depend on their own values. But maybe that's too much for me. Alternatively, split the pliers: give one to the other player and keep one. Then I have one plier (2), chalk (3), nail (2), bucket (1) for a total of 8. The other player gets one plier (2) and whatever's left. Wait, no—if I take one plier, chalk, nail, bucket, that's 2+3+2+1=8. The other player gets the other plier (2). So their total would be 2. That's a big difference. Not good for total.} \\
\\ 
\texttt{Alternatively, maybe split the high{-}value items. Keep chalk (3) and nail (2) for myself, and split the pliers and bucket. So I take chalk, nail, and one plier (total 3+2+2=7). The other player gets the other plier (2) and bucket (1). Their total would be 3. Total combined is 10. That's good. But maybe the other player values the bucket more? Not sure.} \\
\\ 
\texttt{Wait, but the bucket is only} \\
            }
        }
    }
    & & \\ \\

    \theutterance \stepcounter{utterance}  
    & & \multicolumn{4}{p{0.6\linewidth}}{
        \cellcolor[rgb]{0.9,0.9,0.9}{
            \makecell[{{p{\linewidth}}}]{
                \texttt{\tiny{[P1$\langle$GM]}}
                \texttt{<reasoning>} \\
\texttt{Let me try to work through this step by step. The goal is to divide the items so that the sum of both our scores is maximized. Since I don't know the other player's values, I need to make a proposal that's fair and likely to be accepted. Let's start by evaluating the items based on my values.} \\
\\ 
\texttt{My values:} \\
\texttt{{-} Pliers: 2 each (total 4 for both)} \\
\texttt{{-} Chalk: 3} \\
\texttt{{-} Nail: 2} \\
\texttt{{-} Bucket: 1} \\
\\ 
\texttt{Total value I assign: 4 + 3 + 2 + 1 = 10.} \\
\\ 
\texttt{But the other player's values are unknown. However, to maximize the total, I should aim for a division where the sum of both our scores is as high as possible. This might mean balancing the items so that neither of us gets too much or too little.} \\
\\ 
\texttt{Let's list all possible items and think about how to split them. The items are 2 pliers, 1 chalk, 1 nail, 1 bucket. Each player needs to get some subset, and the proposals must be complementary.} \\
\\ 
\texttt{If I take both pliers, that's 4 value for me. Then the other player gets chalk, nail, bucket. Their value for those would depend on their own values. But maybe that's too much for me. Alternatively, split the pliers: give one to the other player and keep one. Then I have one plier (2), chalk (3), nail (2), bucket (1) for a total of 8. The other player gets one plier (2) and whatever's left. Wait, no—if I take one plier, chalk, nail, bucket, that's 2+3+2+1=8. The other player gets the other plier (2). So their total would be 2. That's a big difference. Not good for total.} \\
\\ 
\texttt{Alternatively, maybe split the high{-}value items. Keep chalk (3) and nail (2) for myself, and split the pliers and bucket. So I take chalk, nail, and one plier (total 3+2+2=7). The other player gets the other plier (2) and bucket (1). Their total would be 3. Total combined is 10. That's good. But maybe the other player values the bucket more? Not sure.} \\
\\ 
\texttt{Wait, but the bucket is only} \\
            }
        }
    }
    & \\ \\

    \theutterance \stepcounter{utterance}  
    & \multicolumn{4}{p{0.6\linewidth}}{
        \cellcolor[rgb]{0.9,0.9,0.9}{
            \makecell[{{p{\linewidth}}}]{
                \texttt{\tiny{[P1$\rangle$GM]}}
                \texttt{<reasoning>} \\
\texttt{Let me try to figure out the best way to divide these items. The goal is to maximize the total score for both players. Since I don't know the other player's values, I need to make a proposal that's fair and likely to be accepted. Let's start by evaluating the items based on my values.} \\
\\ 
\texttt{My values:} \\
\texttt{{-} Pliers: 2 each (total 4 for both)} \\
\texttt{{-} Chalk: 3} \\
\texttt{{-} Nail: 2} \\
\texttt{{-} Bucket: 1} \\
\\ 
\texttt{Total value I assign: 4 + 3 + 2 + 1 = 10.} \\
\\ 
\texttt{But the other player's values are unknown. However, to maximize the total, I should aim for a division where the sum of both our scores is as high as possible. This might mean balancing the items so that neither of us gets too much or too little.} \\
\\ 
\texttt{Let's list all possible items and think about how to split them. The items are 2 pliers, 1 chalk, 1 nail, 1 bucket. Each player needs to get some subset, and the proposals must be complementary.} \\
\\ 
\texttt{If I take both pliers, that's 4 value for me. Then the other player gets chalk, nail, bucket. Their value for those would depend on their own values. But maybe that's too much for me. Alternatively, split the pliers: give one to the other player and keep one. Then I have one plier (2), chalk (3), nail (2), bucket (1) for a total of 8. The other player gets one plier (2) and whatever's left. Wait, no—if I take one plier, chalk, nail, bucket, that's 2+3+2+1=8. The other player gets the other plier (2). So their total would be 2. That's a big difference. Not good for total.} \\
\\ 
\texttt{Alternatively, maybe split the high{-}value items. Keep chalk (3) and nail (2) for myself, and split the pliers and bucket. So I take chalk, nail, and one plier (total 3+2+2=7). The other player gets the other plier (2) and bucket (1). Their total would be 3. Total combined is 10. That's good. But maybe the other player values the bucket more? Not sure.} \\
\\ 
\texttt{Wait, but the bucket is only worth 1 to} \\
            }
        }
    }
    & & \\ \\

    \theutterance \stepcounter{utterance}  
    & & & \multicolumn{2}{p{0.3\linewidth}}{
        \cellcolor[rgb]{0.9,0.9,0.9}{
            \makecell[{{p{\linewidth}}}]{
                \texttt{\tiny{[GM$|$GM]}}
                \texttt{continue} \\
            }
        }
    }
    & & \\ \\

    \theutterance \stepcounter{utterance}  
    & & & \multicolumn{2}{p{0.3\linewidth}}{
        \cellcolor[rgb]{0.9,0.9,0.9}{
            \makecell[{{p{\linewidth}}}]{
                \texttt{\tiny{[GM$|$GM]}}
                \texttt{<reasoning>} \\
\texttt{Let me try to figure out the best way to divide these items. The goal is to maximize the total score for both players. Since I don't know the other player's values, I need to make a proposal that's fair and likely to be accepted. Let's start by evaluating the items based on my values.} \\
\\ 
\texttt{My values:} \\
\texttt{{-} Pliers: 2 each (total 4 for both)} \\
\texttt{{-} Chalk: 3} \\
\texttt{{-} Nail: 2} \\
\texttt{{-} Bucket: 1} \\
\\ 
\texttt{Total value I assign: 4 + 3 + 2 + 1 = 10.} \\
\\ 
\texttt{But the other player's values are unknown. However, to maximize the total, I should aim for a division where the sum of both our scores is as high as possible. This might mean balancing the items so that neither of us gets too much or too little.} \\
\\ 
\texttt{Let's list all possible items and think about how to split them. The items are 2 pliers, 1 chalk, 1 nail, 1 bucket. Each player needs to get some subset, and the proposals must be complementary.} \\
\\ 
\texttt{If I take both pliers, that's 4 value for me. Then the other player gets chalk, nail, bucket. Their value for those would depend on their own values. But maybe that's too much for me. Alternatively, split the pliers: give one to the other player and keep one. Then I have one plier (2), chalk (3), nail (2), bucket (1) for a total of 8. The other player gets one plier (2) and whatever's left. Wait, no—if I take one plier, chalk, nail, bucket, that's 2+3+2+1=8. The other player gets the other plier (2). So their total would be 2. That's a big difference. Not good for total.} \\
\\ 
\texttt{Alternatively, maybe split the high{-}value items. Keep chalk (3) and nail (2) for myself, and split the pliers and bucket. So I take chalk, nail, and one plier (total 3+2+2=7). The other player gets the other plier (2) and bucket (1). Their total would be 3. Total combined is 10. That's good. But maybe the other player values the bucket more? Not sure.} \\
\\ 
\texttt{Wait, but the bucket is only worth 1 to} \\
            }
        }
    }
    & & \\ \\

    \theutterance \stepcounter{utterance}  
    & & \multicolumn{4}{p{0.6\linewidth}}{
        \cellcolor[rgb]{0.9,0.9,0.9}{
            \makecell[{{p{\linewidth}}}]{
                \texttt{\tiny{[GM$\rangle$P2]}}
                \texttt{<reasoning>} \\
\texttt{Let me try to figure out the best way to divide these items. The goal is to maximize the total score for both players. Since I don't know the other player's values, I need to make a proposal that's fair and likely to be accepted. Let's start by evaluating the items based on my values.} \\
\\ 
\texttt{My values:} \\
\texttt{{-} Pliers: 2 each (total 4 for both)} \\
\texttt{{-} Chalk: 3} \\
\texttt{{-} Nail: 2} \\
\texttt{{-} Bucket: 1} \\
\\ 
\texttt{Total value I assign: 4 + 3 + 2 + 1 = 10.} \\
\\ 
\texttt{But the other player's values are unknown. However, to maximize the total, I should aim for a division where the sum of both our scores is as high as possible. This might mean balancing the items so that neither of us gets too much or too little.} \\
\\ 
\texttt{Let's list all possible items and think about how to split them. The items are 2 pliers, 1 chalk, 1 nail, 1 bucket. Each player needs to get some subset, and the proposals must be complementary.} \\
\\ 
\texttt{If I take both pliers, that's 4 value for me. Then the other player gets chalk, nail, bucket. Their value for those would depend on their own values. But maybe that's too much for me. Alternatively, split the pliers: give one to the other player and keep one. Then I have one plier (2), chalk (3), nail (2), bucket (1) for a total of 8. The other player gets one plier (2) and whatever's left. Wait, no—if I take one plier, chalk, nail, bucket, that's 2+3+2+1=8. The other player gets the other plier (2). So their total would be 2. That's a big difference. Not good for total.} \\
\\ 
\texttt{Alternatively, maybe split the high{-}value items. Keep chalk (3) and nail (2) for myself, and split the pliers and bucket. So I take chalk, nail, and one plier (total 3+2+2=7). The other player gets the other plier (2) and bucket (1). Their total would be 3. Total combined is 10. That's good. But maybe the other player values the bucket more? Not sure.} \\
\\ 
\texttt{Wait, but the bucket is only worth 1 to} \\
            }
        }
    }
    & \\ \\

    \theutterance \stepcounter{utterance}  
    & & & \multicolumn{4}{p{0.6\linewidth}}{
        \cellcolor[rgb]{0.9,0.9,0.9}{
            \makecell[{{p{\linewidth}}}]{
                \texttt{\tiny{[GM$\langle$P2]}}
                \texttt{<reasoning>} \\
\texttt{Let me try to work through this step by step. The goal is to divide the items so that the sum of both our scores is maximized. Since I don't know the other player's values, I need to make a proposal that's fair and likely to be accepted. Let's start by evaluating the items based on my values.} \\
\\ 
\texttt{My values:} \\
\texttt{{-} Pliers: 2 each (total 4 for both)} \\
\texttt{{-} Chalk: 3} \\
\texttt{{-} Nail: 2} \\
\texttt{{-} Bucket: 1} \\
\\ 
\texttt{Total value I assign: 4 + 3 + 2 + 1 = 10.} \\
\\ 
\texttt{But the other player's values are unknown. However, to maximize the total, I should aim for a division where the sum of both our scores is as high as possible. This might mean balancing the items so that neither of us gets too much or too little.} \\
\\ 
\texttt{Let's list all possible items and think about how to split them. The items are 2 pliers, 1 chalk, 1 nail, 1 bucket. Each player needs to get some subset, and the proposals must be complementary.} \\
\\ 
\texttt{If I take both pliers, that's 4 value for me. Then the other player gets chalk, nail, bucket. Their value for those would depend on their own values. But maybe that's too much for me. Alternatively, split the pliers: give one to the other player and keep one. Then I have one plier (2), chalk (3), nail (2), bucket (1) for a total of 8. The other player gets one plier (2) and whatever's left. Wait, no—if I take one plier, chalk, nail, bucket, that's 2+3+2+1=8. The other player gets the other plier (2). So their total would be 2. That's a big difference. Not good for total.} \\
\\ 
\texttt{Alternatively, maybe split the high{-}value items. Keep chalk (3) and nail (2) for myself, and split the pliers and bucket. So I take chalk, nail, and one plier (total 3+2+2=7). The other player gets the other plier (2) and bucket (1). Their total would be 3. Total combined is 10. That's good. But maybe the other player values the bucket more? Not sure.} \\
\\ 
\texttt{Wait, but the bucket is only} \\
            }
        }
    }
     \\ \\

    \theutterance \stepcounter{utterance}  
    & & & \multicolumn{2}{p{0.3\linewidth}}{
        \cellcolor[rgb]{0.9,0.9,0.9}{
            \makecell[{{p{\linewidth}}}]{
                \texttt{\tiny{[GM$|$GM]}}
                \texttt{continue} \\
            }
        }
    }
    & & \\ \\

    \theutterance \stepcounter{utterance}  
    & & & \multicolumn{2}{p{0.3\linewidth}}{
        \cellcolor[rgb]{0.9,0.9,0.9}{
            \makecell[{{p{\linewidth}}}]{
                \texttt{\tiny{[GM$|$GM]}}
                \texttt{<reasoning>} \\
\texttt{Let me try to work through this step by step. The goal is to divide the items so that the sum of both our scores is maximized. Since I don't know the other player's values, I need to make a proposal that's fair and likely to be accepted. Let's start by evaluating the items based on my values.} \\
\\ 
\texttt{My values:} \\
\texttt{{-} Pliers: 2 each (total 4 for both)} \\
\texttt{{-} Chalk: 3} \\
\texttt{{-} Nail: 2} \\
\texttt{{-} Bucket: 1} \\
\\ 
\texttt{Total value I assign: 4 + 3 + 2 + 1 = 10.} \\
\\ 
\texttt{But the other player's values are unknown. However, to maximize the total, I should aim for a division where the sum of both our scores is as high as possible. This might mean balancing the items so that neither of us gets too much or too little.} \\
\\ 
\texttt{Let's list all possible items and think about how to split them. The items are 2 pliers, 1 chalk, 1 nail, 1 bucket. Each player needs to get some subset, and the proposals must be complementary.} \\
\\ 
\texttt{If I take both pliers, that's 4 value for me. Then the other player gets chalk, nail, bucket. Their value for those would depend on their own values. But maybe that's too much for me. Alternatively, split the pliers: give one to the other player and keep one. Then I have one plier (2), chalk (3), nail (2), bucket (1) for a total of 8. The other player gets one plier (2) and whatever's left. Wait, no—if I take one plier, chalk, nail, bucket, that's 2+3+2+1=8. The other player gets the other plier (2). So their total would be 2. That's a big difference. Not good for total.} \\
\\ 
\texttt{Alternatively, maybe split the high{-}value items. Keep chalk (3) and nail (2) for myself, and split the pliers and bucket. So I take chalk, nail, and one plier (total 3+2+2=7). The other player gets the other plier (2) and bucket (1). Their total would be 3. Total combined is 10. That's good. But maybe the other player values the bucket more? Not sure.} \\
\\ 
\texttt{Wait, but the bucket is only} \\
            }
        }
    }
    & & \\ \\

    \theutterance \stepcounter{utterance}  
    & & \multicolumn{4}{p{0.6\linewidth}}{
        \cellcolor[rgb]{0.9,0.9,0.9}{
            \makecell[{{p{\linewidth}}}]{
                \texttt{\tiny{[P1$\langle$GM]}}
                \texttt{<reasoning>} \\
\texttt{Let me try to work through this step by step. The goal is to divide the items so that the sum of both our scores is maximized. Since I don't know the other player's values, I need to make a proposal that's fair and likely to be accepted. Let's start by evaluating the items based on my values.} \\
\\ 
\texttt{My values:} \\
\texttt{{-} Pliers: 2 each (total 4 for both)} \\
\texttt{{-} Chalk: 3} \\
\texttt{{-} Nail: 2} \\
\texttt{{-} Bucket: 1} \\
\\ 
\texttt{Total value I assign: 4 + 3 + 2 + 1 = 10.} \\
\\ 
\texttt{But the other player's values are unknown. However, to maximize the total, I should aim for a division where the sum of both our scores is as high as possible. This might mean balancing the items so that neither of us gets too much or too little.} \\
\\ 
\texttt{Let's list all possible items and think about how to split them. The items are 2 pliers, 1 chalk, 1 nail, 1 bucket. Each player needs to get some subset, and the proposals must be complementary.} \\
\\ 
\texttt{If I take both pliers, that's 4 value for me. Then the other player gets chalk, nail, bucket. Their value for those would depend on their own values. But maybe that's too much for me. Alternatively, split the pliers: give one to the other player and keep one. Then I have one plier (2), chalk (3), nail (2), bucket (1) for a total of 8. The other player gets one plier (2) and whatever's left. Wait, no—if I take one plier, chalk, nail, bucket, that's 2+3+2+1=8. The other player gets the other plier (2). So their total would be 2. That's a big difference. Not good for total.} \\
\\ 
\texttt{Alternatively, maybe split the high{-}value items. Keep chalk (3) and nail (2) for myself, and split the pliers and bucket. So I take chalk, nail, and one plier (total 3+2+2=7). The other player gets the other plier (2) and bucket (1). Their total would be 3. Total combined is 10. That's good. But maybe the other player values the bucket more? Not sure.} \\
\\ 
\texttt{Wait, but the bucket is only} \\
\\ 
\\ 
\texttt{The time is up. It is now your turn to submit a secret proposal.} \\
            }
        }
    }
    & \\ \\

    \theutterance \stepcounter{utterance}  
    & \multicolumn{4}{p{0.6\linewidth}}{
        \cellcolor[rgb]{0.9,0.9,0.9}{
            \makecell[{{p{\linewidth}}}]{
                \texttt{\tiny{[P1$\rangle$GM]}}
                \texttt{{[}Proposal: 1 plier, chalk, nail{]}} \\
\texttt{</reasoning>} \\
\\ 
\texttt{{[}Proposal: 1 plier, chalk, nail{]}} \\
\\ 
\texttt{ <sep> human:} \\
            }
        }
    }
    & & \\ \\

    \theutterance \stepcounter{utterance}  
    & & & \multicolumn{2}{p{0.3\linewidth}}{
        \cellcolor[rgb]{0.9,0.9,0.9}{
            \makecell[{{p{\linewidth}}}]{
                \texttt{\tiny{[GM$|$GM]}}
                \texttt{proposal does not include number/name pairs} \\
            }
        }
    }
    & & \\ \\

\end{supertabular}
}

\end{document}
